% *======================================================================*
%  Cactus Thorn template for ThornGuide documentation
%  Author: Ian Kelley
%  Date: Sun Jun 02, 2002
%  $Header$
%
%  Thorn documentation in the latex file doc/documentation.tex
%  will be included in ThornGuides built with the Cactus make system.
%  The scripts employed by the make system automatically include
%  pages about variables, parameters and scheduling parsed from the
%  relevant thorn CCL files.
%
%  This template contains guidelines which help to assure that your
%  documentation will be correctly added to ThornGuides. More
%  information is available in the Cactus UsersGuide.
%
%  Guidelines:
%   - Do not change anything before the line
%       % START CACTUS THORNGUIDE",
%     except for filling in the title, author, date, etc. fields.
%        - Each of these fields should only be on ONE line.
%        - Author names should be separated with a \\ or a comma.
%   - You can define your own macros, but they must appear after
%     the START CACTUS THORNGUIDE line, and must not redefine standard
%     latex commands.
%   - To avoid name clashes with other thorns, 'labels', 'citations',
%     'references', and 'image' names should conform to the following
%     convention:
%       ARRANGEMENT_THORN_LABEL
%     For example, an image wave.eps in the arrangement CactusWave and
%     thorn WaveToyC should be renamed to CactusWave_WaveToyC_wave.eps
%   - Graphics should only be included using the graphicx package.
%     More specifically, with the "\includegraphics" command.  Do
%     not specify any graphic file extensions in your .tex file. This
%     will allow us to create a PDF version of the ThornGuide
%     via pdflatex.
%   - References should be included with the latex "\bibitem" command.
%   - Use \begin{abstract}...\end{abstract} instead of \abstract{...}
%   - Do not use \appendix, instead include any appendices you need as
%     standard sections.
%   - For the benefit of our Perl scripts, and for future extensions,
%     please use simple latex.
%
% *======================================================================*
%
% Example of including a graphic image:
%    \begin{figure}[ht]
% 	\begin{center}
%    	   \includegraphics[width=6cm]{MyArrangement_MyThorn_MyFigure}
% 	\end{center}
% 	\caption{Illustration of this and that}
% 	\label{MyArrangement_MyThorn_MyLabel}
%    \end{figure}
%
% Example of using a label:
%   \label{MyArrangement_MyThorn_MyLabel}
%
% Example of a citation:
%    \cite{MyArrangement_MyThorn_Author99}
%
% Example of including a reference
%   \bibitem{MyArrangement_MyThorn_Author99}
%   {J. Author, {\em The Title of the Book, Journal, or periodical}, 1 (1999),
%   1--16. {\tt http://www.nowhere.com/}}
%
% *======================================================================*

% If you are using CVS use this line to give version information
% $Header$

\documentclass{article}

% Use the Cactus ThornGuide style file
% (Automatically used from Cactus distribution, if you have a
%  thorn without the Cactus Flesh download this from the Cactus
%  homepage at www.cactuscode.org)
\usepackage{../../../../doc/latex/cactus}

\begin{document}

% The author of the documentation
\author{Erik Schnetter \textless schnetter@gmail.com\textgreater}

% The title of the document (not necessarily the name of the Thorn)
\title{SystemTopology}

% the date your document was last changed, if your document is in CVS,
% please use:
%    \date{$ $Date$ $}
\date{July 13 2026}

\maketitle

% Do not delete next line
% START CACTUS THORNGUIDE

% Add all definitions used in this documentation here
%   \def\mydef etc

% Add an abstract for this thorn's documentation
\begin{abstract}
  At startup, report the MPI process/host mapping and the node's
  hardware topology (cores, PUs, caches, NUMA memory) as discovered by
  \emph{hwloc}, optionally bind OpenMP threads (and their memory
  allocations) to specific processing units, and expose the discovered
  topology to other thorns through aliased functions.
\end{abstract}

\section{Introduction}

This thorn runs once, early during startup (scheduled at
\texttt{AT startup AFTER hwloc\_startup BEFORE Driver\_Startup}), and
performs three tasks:
\begin{enumerate}
\item It determines how MPI processes are distributed across physical
  hosts (which processes share a host, and in what order), and prints
  this mapping.
\item It queries \emph{hwloc} for the local node's hardware topology
  (cores, hardware threads / PUs, cache hierarchy, NUMA nodes) and
  prints it, together with a report of which \emph{hwloc} features
  (binding, discovery) are actually supported on this system.
\item Optionally, it binds each OpenMP thread of each MPI process on a
  host to a specific processing unit (PU), and binds that thread's
  memory allocations to the same PU's local memory.
\end{enumerate}
The topology information gathered here is not only diagnostic output:
it is also captured and made available to other thorns (see
section~\ref{sec:interaction}), which use it to make cache- and
NUMA-aware decisions about loop tiling and memory allocation.

This thorn requires the \texttt{hwloc} and \texttt{MPI} capabilities
(see \texttt{configuration.ccl}).

\section{Node and MPI Topology Discovery}

Discovery proceeds in two independent steps.

\paragraph{MPI process/host mapping.} Each MPI process determines its
hostname (via \texttt{MPI\_Get\_processor\_name}) and these are
gathered on the global root process, which groups processes into
hosts by matching hostnames, and computes, for every process: which
host it runs on, how many processes run on that host, and this
process's position (rank) among them. This mapping is then scattered
back to every process. Each consumer of this mapping (see below) that
also needs per-thread information -- e.g. to lay out OpenMP threads
across a host's processes -- combines it locally with its own
\texttt{omp\_get\_max\_threads()} (or \texttt{omp\_get\_num\_threads()}),
assuming every rank on a host runs the same number of OpenMP threads
(see the ``all nodes identical'' assumption below); no further MPI
communication is needed for this.

\paragraph{Hardware topology.} \emph{hwloc} is used to enumerate cores,
PUs (hardware threads), the cache hierarchy, and NUMA/machine memory
nodes. From this, the thorn derives:
\begin{itemize}
\item the number of PUs per core (the hardware SMT-thread count), and
\item the number of SMT threads actually used, i.e.
  $\lceil (\text{total OpenMP threads on this host}) /
  (\text{number of cores}) \rceil$.
\end{itemize}
If the process is already confined to a subset of the node's PUs --
for example by a job scheduler's cpuset/cgroup, or by the MPI
launcher's own affinity settings (\texttt{srun --cpu-bind},
\texttt{mpirun --bind-to}, etc.) -- the \emph{hwloc} topology is
restricted to that subset before any of the above is computed, so that
core/PU counts, bindings, and cache/memory sizes all reflect only the
resources actually available to this process.

On a heterogeneous (hybrid) CPU -- one with distinct ``performance''
and ``efficiency'' core kinds, such as recent Intel P-core/E-core
parts or ARM big.LITTLE designs -- the topology can be further
restricted to the highest-efficiency core kind only, using
\emph{hwloc}'s \texttt{cpukinds} API (see the
\texttt{CACTUS\_USE\_PERFORMANCE\_CORES\_ONLY} environment variable
below). This restriction is applied on top of, and intersected with,
any scheduler/launcher restriction already in effect above.

This thorn assumes that \textbf{all nodes used in a run are
identical}. Consequently, only the global root process (MPI rank~0)
actually queries \emph{hwloc}'s \texttt{cpukinds} API; the outcome
(whether performance cores could be reliably identified, and if so,
their PU set) is broadcast via MPI to every other rank, which then
each locally intersects it with their own scheduler/launcher
restriction from above. This avoids every process on every node
redundantly repeating the same query, but means that on a genuinely
heterogeneous cluster (where some nodes are hybrid CPUs and others are
not, or different nodes have different performance-core layouts), the
root node's detection would be incorrectly applied to every other
node.

\section{Thread and Memory Binding}

If enabled (see the \texttt{set\_thread\_bindings} parameter), each
OpenMP thread is bound to exactly one PU using the following layout
algorithm:
\begin{itemize}
\item If \texttt{thread\_layout = "dense"}, threads are packed onto as
  few cores as possible (filling all SMT threads of a core before
  moving to the next core).
\item If \texttt{thread\_layout = "loose"}, threads are spread evenly
  across all available cores.
\end{itemize}
Threads are numbered consecutively across all MPI processes on a host
(using the per-rank thread offsets described above), then mapped onto
cores and PUs; if there are more threads than PUs, cores/PUs are
oversubscribed rather than left unbound. When a PU binding succeeds,
the thorn additionally binds that thread's future memory allocations
to the same PU's local memory, if \emph{hwloc} reports support for
this on the current system.

The memory binding \emph{policy} is treated as a property of the
machine, not of the simulation, and is therefore controlled by the
environment variable \texttt{CACTUS\_MEMBIND\_POLICY} rather than a
Cactus parameter (following the convention of e.g.
\texttt{CACTUS\_NUM\_THREADS}):
\begin{itemize}
\item \texttt{"bind"} (the default, used if the variable is unset):
  strictly bind allocations to the thread's local NUMA node; if that
  node is full, allocation fails rather than silently using a
  different node.

  \emph{Use this when} each thread's memory footprint is dominated by
  data it privately owns and repeatedly accesses -- the common case for
  a Carpet/CarpetX run, where each OpenMP thread mostly reads and
  writes grid-function data within its own loop tiles (see
  section~\ref{sec:interaction}) -- and per-thread memory usage is
  reasonably balanced across the node, so no single NUMA node is likely
  to run out. This is the right default for most production runs: it
  gives deterministic, guaranteed locality rather than relying on the
  OS's first-touch heuristic, which can be defeated by memory pools,
  pre-faulted pages, or a cluster-wide NUMA policy set by an
  administrator or container runtime.
\item \texttt{"firsttouch"}: memory is placed on the NUMA node of
  whichever thread first touches it (this approximates, and makes
  explicit, the usual OS default behavior with no binding at all).

  \emph{Use this when} you want the same intended locality as
  \texttt{"bind"}, but are not confident that per-thread memory usage
  is well balanced (e.g. a run with a very uneven domain decomposition
  or adaptive mesh refinement concentrating much more grid data on some
  threads/cores than others). Unlike \texttt{"bind"}, a full local node
  degrades performance instead of aborting the run outright. The
  tradeoff is that first-touch placement depends on which thread
  genuinely touches a page \emph{first}; if grid-function memory is
  pre-allocated or pre-zeroed by a single setup thread before the
  worker threads that will actually use it begin their loop tiles, the
  intended locality can be lost regardless of this setting.
\item \texttt{"interleave"}: memory pages are spread round-robin
  across NUMA nodes; useful for large structures deliberately shared
  across sockets.

  \emph{Use this when} a particular allocation is genuinely read (or
  written) roughly equally by threads across multiple NUMA nodes, so no
  single ``local'' placement is meaningful -- for example, a large
  equation-of-state or interpolation table shared by all threads on a
  host. Interleaving spreads the aggregate memory-bandwidth demand
  across all memory controllers on the node instead of concentrating it
  on one, which would otherwise become a bottleneck (and would not
  actually be ``local'' to most of the threads hammering it anyway).
  This is a niche choice, not appropriate for the bulk of
  per-thread grid-function storage.
\item \texttt{"none"}: do not set any memory binding at all (CPU
  binding, if enabled, is unaffected).

  \emph{Use this when} diagnosing a performance or stability issue and
  you want to isolate whether CPU pinning or memory binding is
  responsible, when an external tool (\texttt{numactl}, a container
  runtime, or a job scheduler) already manages memory placement for the
  job and you do not want this thorn's choice to conflict with it, or
  on a system where \emph{hwloc}'s memory-binding support is unreliable
  even though CPU-binding support works.
\end{itemize}
If \emph{hwloc} reports that the requested policy is not supported on
this system, the thorn logs a warning and does not set memory bindings
for this run. An unrecognized value aborts with an error.

\section{Using This Thorn}

\subsection{Obtaining This Thorn}

This thorn is part of the \texttt{CactusUtils} arrangement and is
obtained together with the rest of the Cactus Einstein Toolkit
checkout.

\subsection{Basic Usage}

Simply activate the thorn; it requires no scheduling from other
thorns and runs automatically once at startup. Its behaviour is
controlled by two steerable parameters:
\begin{itemize}
\item \texttt{set\_thread\_bindings} (\texttt{"yes"} / \texttt{"no"} /
  \texttt{"auto"}, default \texttt{"auto"}): whether to set thread-CPU
  affinity at all. \texttt{"auto"} is the right choice for most runs.
  Use \texttt{"no"} when the job launcher's own affinity flags
  (\texttt{srun --cpu-bind}, \texttt{mpirun --bind-to}) already do what
  you want and this thorn should not touch it, or when isolating
  whether CPU pinning itself is responsible for an observed
  performance or correctness issue. Use \texttt{"yes"} to force this
  thorn's own layout even in cases where it would otherwise detect an
  existing restriction and back off (see ``\texttt{"auto"} binding''
  below) -- only appropriate if you know that restriction is still
  compatible with this thorn's layout arithmetic (e.g. a plain,
  contiguous per-rank slice of the node).
\item \texttt{thread\_layout} (\texttt{"dense"} / \texttt{"loose"},
  default \texttt{"dense"}): how threads are laid out over cores when
  bindings are set. This only matters when a host's total thread count
  (summed over all ranks on it) is \emph{less} than its core count --
  otherwise dense and loose produce the same layout. \texttt{"dense"}
  packs threads onto as few cores (and ideally as few NUMA nodes) as
  possible; choose this when several ranks share a host and each
  should stay confined to its own socket (avoiding cross-socket memory
  traffic for that rank's own data), or when leaving cores fully idle
  is desirable (e.g. for power/thermal reasons, or to leave headroom
  for co-located non-Cactus processes). \texttt{"loose"} spreads
  threads across all available cores/sockets; choose this for a single
  under-subscribed rank that should exploit as much aggregate memory
  bandwidth and cache capacity as possible, accepting the resulting
  cross-socket NUMA effects within that rank's own address space.
\end{itemize}
Two further settings are controlled by environment variables rather
than parameters, since they describe properties of the machine rather
than of the simulation (see the sections referenced below for details):
\begin{itemize}
\item \texttt{CACTUS\_USE\_PERFORMANCE\_CORES\_ONLY} (\texttt{"yes"} /
  \texttt{"no"} / \texttt{"auto"}, default \texttt{"auto"}): whether to
  restrict this thorn's view of the node -- core/PU counts, cache
  info, and thread bindings -- to performance cores only, on a
  heterogeneous CPU. \texttt{"auto"} is the right choice for most runs:
  it uses performance cores when it can identify them and degrades
  gracefully otherwise. Use \texttt{"yes"} when you need a hard
  guarantee that efficiency cores are never silently included -- e.g.
  for reproducible benchmarking, or on a shared workstation/laptop
  where mixing in much slower efficiency cores would badly skew
  load-balancing assumptions -- and would rather the run abort at
  startup than proceed with an uncertain core set. Use \texttt{"no"}
  when you deliberately want every core, efficiency cores included,
  such as for a lower-priority or embarrassingly parallel run where
  raw throughput matters more than uniform per-thread performance.
\item \texttt{CACTUS\_MEMBIND\_POLICY} (\texttt{"bind"} /
  \texttt{"firsttouch"} / \texttt{"interleave"} / \texttt{"none"},
  default \texttt{"bind"}): the NUMA memory binding policy applied
  alongside CPU bindings.
\end{itemize}

\subsection{Special Behaviour}

\begin{itemize}
\item \textbf{\texttt{"auto"} binding.} In \texttt{"auto"} mode, the
  thorn only sets thread bindings if this process was \emph{not}
  already confined to a subset of the node by an external mechanism
  (job scheduler, cgroup, or launcher affinity flags). If such a
  restriction is detected, \texttt{"auto"} behaves like \texttt{"no"},
  since the thorn's own whole-node layout arithmetic does not apply to
  a partial-node allocation and would only fight whatever affinity has
  already been set. Use \texttt{set\_thread\_bindings = "yes"} to
  force the thorn's own layout regardless.
\item \textbf{Uniform thread counts.} Since all nodes are assumed
  identical, every rank on a host is assumed to run the same number of
  OpenMP threads; this count is used directly (no MPI communication is
  needed to determine it) to compute each rank's thread offset and the
  host's total thread count.
\item \textbf{Performance-core restriction.} With
  \texttt{CACTUS\_USE\_PERFORMANCE\_CORES\_ONLY=auto} (the default), the
  thorn restricts its view of the node to performance cores only if it
  can reliably identify them via \emph{hwloc}'s \texttt{cpukinds} API
  (requires \emph{hwloc} $\geq$~2.4.0 and operating-system/hardware
  support for reporting core efficiency); otherwise it quietly falls
  back to using all cores. With \texttt{=yes}, the same restriction is
  mandatory: if performance cores cannot be reliably identified
  (unsupported \emph{hwloc} version, a homogeneous CPU, an unknown
  efficiency ranking, or a scheduler/launcher restriction that excludes
  all performance cores), the run aborts with an error rather than
  silently using efficiency cores. \texttt{=no} disables this
  restriction entirely.
\item \textbf{Binding report.} The ``Thread CPU bindings'' report is
  collected and printed by global MPI rank~0, which is always process~0
  on host~0 (host~0 is defined as whichever host rank~0 happens to run
  on, and rank~0 -- being the smallest possible rank -- is always the
  smallest, i.e. ``process~0'', on whatever host it shares).
\item \textbf{Missing hwloc support.} If \emph{hwloc} reports that it
  cannot query or set CPU/memory bindings on this system (see the
  ``Topology support'' report), the corresponding step is skipped with
  an informational message rather than failing.
\end{itemize}

\subsection{Interaction With Other Thorns}
\label{sec:interaction}

This thorn provides four aliased functions (see
\texttt{interface.ccl}) that other thorns use to make cache- and
NUMA-aware decisions:
\begin{itemize}
\item \texttt{GetNumSMTThreads} / \texttt{GetMaxSMTThreads} -- the
  number of SMT threads in use and the hardware maximum. Used by
  \texttt{Carpet/LoopControl} to size the innermost loop tiles so that
  hyperthread siblings do not needlessly duplicate cache-line work.
\item \texttt{GetCacheInfo1} -- sizes, line sizes, and strides of each
  cache level, plus the size of local and global (NUMA) memory. Used by
  \texttt{Carpet/LoopControl} and \texttt{Carpet/CarpetLib} to choose
  cache-blocking tile sizes and allocation padding/alignment, and by
  \texttt{CactusUtils/MemSpeed} to select the array sizes used for its
  cache/memory benchmarks.
\item \texttt{GetMPIProcessInfo} -- the MPI process/host mapping
  computed above. Used by \texttt{CactusUtils/MemSpeed} to restrict its
  (expensive) benchmarks to one process per host.
\end{itemize}
All of these thorns query
\texttt{CCTK\_IsFunctionAliased} before calling these functions and
fall back to conservative defaults (e.g. one SMT thread, unknown cache
sizes) if this thorn is not active, so activating this thorn is
optional but recommended whenever these other thorns are in use.

\subsection{Examples}

Illustrative (abbreviated) startup output on a two-socket, multi-core
node:
\begin{verbatim}
INFO (SystemTopology): MPI process-to-host mapping:
This is MPI process 0 of 4
MPI hosts:
  0: node042
This MPI process runs on host 0 of 1
On this host, this is MPI process 0 of 4
INFO (SystemTopology): Topology support:
Discovery support:
  discovery->pu                           : yes
CPU binding support:
  cpubind->set_thisthread_cpubind         : yes
  ...
INFO (SystemTopology): Hardware objects in this node:
Machine L#0: (P#0)
  Package L#0: (P#0)
    Core L#0: (P#0)
      PU L#0: (P#0)
      PU L#1: (P#16)
  ...
INFO (SystemTopology): Setting thread CPU bindings:
INFO (SystemTopology): Extracting CPU/cache/memory properties:
  There are 2 PUs per core (aka hardware SMT threads)
  There are 1 threads per core (aka SMT threads used)
  Cache D1 cache has type "data" depth 1
    size 32768 linesize 64 associativity 8 stride 4096, for 2 PUs
  ...
\end{verbatim}

On a heterogeneous (hybrid) CPU, with
\texttt{CACTUS\_USE\_PERFORMANCE\_CORES\_ONLY} unset, \texttt{auto}, or
\texttt{yes}, an additional message appears before the topology is
reported, and all subsequent core/PU/cache counts reflect performance
cores only:
\begin{verbatim}
INFO (SystemTopology): Detected a heterogeneous (hybrid) CPU;
restricting to performance cores only
\end{verbatim}

\subsection{Support and Feedback}

Contact the maintainer listed in the \texttt{README} file, or use the
Cactus/Einstein Toolkit support channels.

\section{Caveats}

\begin{itemize}
\item The quality of the ``dense''/``loose'' layout depends on
  \emph{hwloc}'s logical ordering of cores actually reflecting NUMA
  locality (adjacent logical indices on the same socket), which holds
  for essentially all common systems but is not guaranteed by
  \emph{hwloc} itself.
\item \texttt{GetCacheInfo1} may report a cache line size of zero if
  \emph{hwloc} cannot determine it (seen inside some virtual machines,
  containers, or on less common architectures); callers should treat
  this as ``unknown'' rather than a literal zero-byte line.
\item Distance-matrix reporting requires \emph{hwloc} version~2.0 or
  later; on older \emph{hwloc} installations this part of the report is
  omitted.
\item Performance-core detection requires \emph{hwloc} version~2.4.0 or
  later, \emph{and} an operating system/hardware combination that
  exposes core efficiency information to \emph{hwloc} (e.g. via
  \texttt{cpu\_capacity} in Linux \texttt{sysfs}, or natively on
  Windows~10 and macOS). On other systems, \texttt{"auto"} simply does
  not restrict, while \texttt{"yes"} aborts with an error.
\item Performance-core detection assumes all nodes in the run are
  identical: it is only performed once, on the global root process,
  and the result is broadcast to all other ranks rather than
  independently rediscovered on every node. On a cluster with a
  genuinely mixed node population, this would incorrectly apply the
  root node's hybrid-CPU layout (or lack thereof) to every other node.
\item Thread-layout arithmetic assumes every rank on a host runs the
  same number of OpenMP threads as every other rank on that host (this
  is not communicated or verified across ranks). A run that
  deliberately assigns different thread counts to different ranks on
  the same host will get an incorrect thread offset/layout computation.
\end{itemize}

\section{History}

\subsection{Thorn Source Code}

The source code lives in \texttt{CactusUtils/SystemTopology} in the
Cactus Einstein Toolkit checkout.

\subsection{Thorn Documentation}

This document was last substantially revised in July 2026, to describe
the handling of externally-restricted (job-scheduler/cgroup) CPU sets,
NUMA memory binding, distance-matrix reporting, and restricting to
performance cores only on heterogeneous (hybrid) CPUs -- the latter,
along with the assumption of uniform per-rank thread counts, relying on
all nodes used in a run being identical.

\subsection{Acknowledgements}

Written and maintained by Erik Schnetter.


\begin{thebibliography}{9}

\end{thebibliography}

% Do not delete next line
% END CACTUS THORNGUIDE

\end{document}
